\documentclass[12pt]{article}
\usepackage{graphicx}
\usepackage{hyperref}
\usepackage{enumitem}
\usepackage{listings}
\usepackage{geometry}

% Adjust margins to give more space
\geometry{a4paper, margin=0.75in}

\title{Weekly Report: LLM + RAG Model Project}
\author{Edward Yang}
\date{\today}

\begin{document}

\maketitle

\section*{Week Overview: 280425 - 050525}
\begin{itemize}
    \item \textbf{Project Name:} An evaluation of Retrieval-Augmented Generation (RAG) Frameworks for Disease Diagnosis
    \item \textbf{Project Description:} This project aims to evaluate the performance of Retrieval Augmented Generation (RAG) frameworks compared to generic LLMs for medical disease diagnosis. Specifically, we use the ICD-10 Disease Dataset, which contains a hierarchical structure of diseases and their corresponding codes to evaluate the accuracy of the RAG model in generating responses to medical queries.
    \item \textbf{Overview:}
    \begin{itemize}[leftmargin=*]
        \item \textbf{Goal 1:} Create Document preprocessing functions and Embedding Model for the Disease Dataset
        \item \textbf{Goal 2:} Use vector database (ChromaDB) to store the embeddings
        \item \textbf{Goal 3:} Retrieval of samples using similiarity search, compared with query
        \item \textbf{Goal 4:} Integrate the RAG model with the LLM for generating responses
        \item \textbf{Goal 5:} Test the model with sample queries
        \item \textbf{Goal 6:} Compare the model's output accuracy against the Non-RAG model (LLM-only). This can be done by creating a fine-tuning model to evaluate diagnostic accuracy of the two models, trained on the hierarchical ICD-10 Disease Dataset, where partial matches at different levels of specificity have different clinical value (similar to the method in the Medfound paper).
        \item \textbf{Goal 7:} Create a report on the model's performance and potential improvements
    \end{itemize}
\end{itemize}

\section*{Progress}
\begin{itemize}
    \item \textbf{Completed Tasks:}
    \begin{itemize}[leftmargin=*]
        \item \textbf{Task 1:} Created data\_loader.py and other tools for loading and preprocessing the dataset.
        \item \textbf{Task 2:} Started on the embedding model, using the Sentence-BERT model.
        \item \textbf{Task 3:} Used ChromaDB as a vector database for storing the embeddings.
        \item \textbf{Task 4:} Implemented similiarity search using ChromaDB query function to retrieve top k most relevant documents based on the embedded query.
        \item \textbf{Task 5:} Integrated Google's Gemma Huggingface model with varied parameter sizes (e.g., 2B, 6B) for generating responses.
    \end{itemize}
    \item \textbf{Ongoing Tasks:}
    \begin{itemize}[leftmargin=*]
        \item \textbf{Task 1:} Create Query template for the RAG model to generate responses.
        \item \textbf{Task 2:} Create the fine tuning model to evaluate the RAG model's accuracy against the LLM-only model.
        \item \textbf{Task 3:} Create a report on the model's performance and potential improvements.
    \end{itemize}
\end{itemize}


\end{document}
